\documentclass[a4paper,11pt]{article}

\usepackage{revy}
\usepackage[utf8]{inputenc}
\usepackage[T1]{fontenc}
\usepackage[danish]{babel}


\revyname{MatematikRevyen}
\revyyear{2026}
% HUSK AT OPDATERE VERSIONSNUMMER
\version{1}
\eta{$4$ minutter}
\status{Færdig}

\title{Mine rus kan tæve dine rus!}
\author{Julius '24, Frederikke '24, Louise '23, Christian (Øko) '24}

\begin{document}
\maketitle

\begin{roles}
\role{M1} Mentor for "rødt"\ mentorhold
\role{M2} Mentor for "blåt"\ mentorhold
\role{R1} Rus på rødt hold der har sort bælte i karate
\role{R2} Rus på blåt hold der har været i Den Kongelige Livgarde
\role{R3} Rus på rødt hold der kan løse en Rubik's Cube på under 7 sekunder
\role{R4} Rus på blåt hold der er IM i skak
\role{R5} Rus på rødt hold der er Færøernes hurtigste sprinter
\role{R6} Rus på blåt hold der er fucking god til matematik
\end{roles}

\begin{props}
\prop{Kager} To veludførte kager
\prop{Dør} En dør til et A-lokale
\prop{Madras/boksering} Underlag til wrestling
\prop{Diverse rekvisitter til slåskamp}
\end{props}


\begin{sketch}

\scene{M1 står ved døren på midten af scenen. M2 kommer ind fra siden. De har hver en stor kage i hænderne} 

\says{M1} Hold da op en kage, du har med. Du skal også holde mentormøde?

\says{M2} Ja, der er jo torsdag. Men der er jo, som sædvanlig, ikke blevet booket lokaler til os. Og der er undervisning i alle lokalerne undtagen A101, selvom der er frokostpause. Så det snupper vi, når MatIntro er slut. Hvor skal I være?

\says{M1} Øh.. JEG havde ligesom tænkt mig at tage A101. Du må finde et andet sted.

\says{M2} Men der er ikke andre steder.

\says{M1} Du kan ikke bare tage mit lokale! Jeg har allerede sagt til mine rus, at det er her vi er!

\says{M2} Det er sgu da ikke dit lokale, hvis du ikke har booket det!

\says{M1} Det er det da ihvertfald! Jeg var her først!

\says{M2} Jeg var her sgu da ligeså først! ... Jeg tror bare jeg går ind i lokalet NU og hører det sidste af øvelsestimen. Jeg kunne alligevel godt bruge en genopfrisker af, hvordan man differentierer.

\scene{M2 går mod døren. M1 løber efter og får stoppet M2.}

\says{M1} Det kan du ik!

\says{M2} Hvorfor ik'?

\says{M1} Fint så. Vi erobrer det bare tilbage. Mine rus kan tæve dine rus!

\scene{Bandet optrapper stemningen med suspenseful musik, følgende dialog, sker i takt til musikken}

\says{M2} Ah hva'?

\says{M1} Jeg sagde mine rus kan tæve dine rus!

\says{M2} NEJ!... MINE rus kan tæve DINE rus!

\says{M1} Jeg har en rus, der har det sorte bælte i karate!

\scene{R1 kommer ud af A101 med gi og sortbælte på og stiller sig bag M1 og laver kiai}

\says{M2} Jeg har en rus der har været i Den Kongelige Livgarde!

\scene{R2 kommer ud med bjørneskindshue, marcherende og stiller sig bag M2}

\says{M1} Jeg har en rus, der kan løse en Rubik's Cube på under 7 sekunder!

\scene{R3 kommer ud med en Rubik's Cube og stiller sig bag M1}

\says{M2} Jeg har en rus der er international mester i skak!

\scene{R4 kommer ud med et skakbræt og stiller sig bag M2}

\says{M1} Jeg har en rus der er Færøernes hurtigste sprinter!

\scene{R5 kommer ud i løbetøj og hurtigbriller og stiller sig bag M1}

\says{M2} Jeg har en rus der er fucking god til matematik!

\scene{R6 kommer ud og stiller sig bag M2. Madrassen bliver båret ind af ninjaer.}

\says{M1} Mine rus kan tæve dine rus! \act{De "røde"\ rus stiller sig op på madrassen, kampklar}

\says{M2} Nej mine rus kan tæve dine rus! \act{De "blå"\ rus stiller sig op på madrassen ansigt til ansigt med de røde}

\scene{Musikken når sit klimaks og holder.}

\says{AV} 3... 2... 1... FIGHT!

\scene{Beatet dropper og bandet spiller hurtig og hård musik. Scenen fyldes med røg og farvet strobelys, så det er svært at se præcis hvad der foregår. Rus slås pro wrestling style. Bruger enhver rekvisit som våben: Tomme GT-dåser, stole, skakbrikker, mentorkager, døre, evt. rekvisitter fra tidligere akter, etc. Mentorerne står på hver deres side af madrassen og propper et styk kage i munden på slåede rus, så de bliver kampklar igen. Kampen varer cirka et minut.}

\scene{Musikken slutter, og lyset bliver normalt igen. På gulvet ligger alle seks rus besejret, ukampdygtige.}

\says{M1} Hvad øh? Vi kan også bare holde mentormøde sammen?

\scene{Lys ned.}

\end{sketch}

\scene{Note: Denne sketch blev skrevet sidste år og er baseret på ægte bedrifter rus fra årgang '25 har opnået. Måske den kan skrives om til at være baseret på ting personer fra årgang '26 har opnået. (Jeg kender dem ikke. Er på udveksling.)}

\end{document}
